\raggedbottom

\usepackage{color, xifthen, amsthm, amssymb, amsmath, tikz, pgfplots}
\pgfplotsset{compat=1.18}
\usetikzlibrary{backgrounds}
\usepackage[letterpaper, margin=1in]{geometry}
\usepackage{multicol}
\usepackage{fourier}  %don't use txfonts b/c the sf doesn't look good

\usepackage{enumitem}
\setlist{topsep=-\parskip/2}


\setlength{\parindent}{0pt}
\setlength{\parskip}{\baselineskip}


\usepackage{mdframed} % makes better frames around examples (eg pagebreak allowed)
  % from mdframed package. used for examples
\newmdenv[
  skipabove=10pt,%\topsep,
  skipbelow=10pt,%\topsep
  ]{allrules}

\usepackage{tcolorbox} % for colored boxes


\usepackage{subfigure}
\usepackage{hyperref}
\hypersetup{colorlinks, urlcolor=red, linkcolor=.} %color urls, but not interal refs
\usepackage{graphicx, floatflt}
\graphicspath{{./figs/}{../figs/}}
\DeclareGraphicsExtensions{.png,.jpg,.pdf}

\usepackage{wrapfig}  %useful for wrapping text around pictures or tables
		      %though it doesn't do a great job inside my example boxes


\definecolor{defbod}{rgb}{.8,.8,.8}     \definecolor{defhead}{rgb}{.7,.1,.1}
\definecolor{thmbod}{rgb}{.8,.8,.8}     \definecolor{thmhead}{rgb}{.1,.1,.9}
\definecolor{factbod}{rgb}{1,1,1}    \definecolor{facthead}{rgb}{.1,.7,.1}
\definecolor{exbod}{rgb}{1,1,1}         \definecolor{exhead}{rgb}{0,.3,0}
\definecolor{soltext}{rgb}{.2,.2,.5}
\definecolor{notebod}{rgb}{.9,.9,.9}    \definecolor{notehead}{rgb}{.1,.1,.7}

\newcommand{\defn}[1]{
 \noindent \colorbox{defbod}{\begin{minipage}{.95\linewidth}
 \hspace{-12pt}\colorbox{defhead}{\sf\color{white} Definition}
 #1
 \end{minipage} } \medskip
 }
\newcommand{\dt}[1]{\textcolor{defhead}{\bfseries #1}}

\newcommand{\thrm}[1]{
 \noindent \colorbox{thmbod}{\begin{minipage}{.95\linewidth}
 \hspace{-12pt}\colorbox{thmhead}{\sf\color{white}Theorem}
 #1
 \end{minipage} } \medskip
 }
\newcommand{\coro}[1]{
 \noindent \colorbox{thmbod}{\begin{minipage}{.95\linewidth}
 \hspace{-12pt}\colorbox{thmhead}{\sf\color{white}Corollary}
 #1
 \end{minipage} } \medskip
 }
\newcommand{\fact}[1]{  %%%% !! note this different from calcnotes
 \noindent \begin{allrules}
 #1
 \end{allrules} \medskip
 }

\newcommand{\note}[1]{
	\noindent \colorbox{notebod}{\begin{minipage}{.95\textwidth}
			\hspace{-12pt}\colorbox{notehead}{\sf\color{white}NOTE}
			#1
	\end{minipage} } \medskip
}


% making the example thing with a line down the left.
\newmdenv[
rightline=false,
bottomline=false,
skipabove=0pt,%\topsep,
skipbelow=0pt,%\topsep
]{siderules}

\newcommand{\ex}[1]{
  \noindent \begin{siderules}
    \textcolor{exhead}{\sf Example}
    #1
  \end{siderules} 
  \medskip
}




%%%%%%%%%% set it up so that these are not displayed when 
%%%%%%%%%% printing notes for students?  or for overhead?
\newcommand{\exsol}[1]{
  {\color{soltext} #1}
}

\newcommand{\mynonumsect}[2][]{
 \cleardoublepage
 \section*{#1\hskip16pt #2}
 \addcontentsline{toc}{section}{#1 #2}
}

\newcommand{\myhead}[1]{
\par%\vskip\baselineskip
\noindent
{\large\bf #1}
}

%%%%%%%%%%%%%%%%%%%%%%%%
% symbol shortcuts

\newcommand{\nice}{\displaystyle}
\newcommand{\eqq}{\ensuremath \stackrel{?}{=}}
\newcommand{\R}{\ensuremath{\mathbb{R}}}




% %%%%%%%%%%%%%%%%%%%%%%%%
% %
% % Creating handouts vs notes
% %% If you use the line below before loading this file then it will be set up for
% %% lecture notes.  Default is for handouts
% %%
% %% \newcommand{\lectureNoteMode}{}
% %%
% %% do not uncomment here, just add the line to the file where you input this one.
% \ifthenelse{\not\isundefined\lectureNoteMode}
% { % lecture notes 
%   \newcommand{\Hvskip}[1]{}
%   \newcommand{\Hvfill}{}
%   \newcommand{\Hpbreak}{}  % pagebreaks for handouts
%   \newcommand{\Npbreak}{\vfill\pagebreak} % pagebreaks for printable notes
%   \newcommand{\Nonly}[1]{#1} % shows in notes, not in handouts
%   \newcommand{\mycleardoublepage}{\cleardoublepage}

%   \newcommand{\ex}[1]{
%    \noindent \begin{allrules}
%     \textcolor{exhead}{\sf Example}\par
%     #1
%     \end{allrules} \medskip
%   }

% } 
% { 
% %% Handout mode 
% %% this section sets it up with smaller margins, and blank spots for writing solns to examples
%   %%%%%%%%%%%%%%%%%%%%%%%%%%%%%%%%%%%%%%%%%%
%   %change frame around examples
%   % from mdframed package
%   \newmdenv[
%   rightline=false,
%   bottomline=false,
%   skipabove=0pt,%\topsep,
%   skipbelow=0pt,%\topsep
%   ]{siderules}
  
%   \newcommand{\ex}[1]{
%     \noindent \begin{siderules}
%       \textcolor{exhead}{\sf Example}
%       #1
%     \end{siderules} 
%     \medskip
%   }
  
%   %
%   %%%%%%%%%%%%%%%%%%%%%%%%%%%%%%%%%%%%%%%%%%
  
%   % don't show solutions for examples
%   \renewcommand{\exsol}[1]{}
  
%   %used for making blank space in handouts. eg: \vs2
%   \newcommand{\Hvskip}[1]{\vskip#1in\null}
%   \newcommand{\Hvfill}{\vfill} %vfill only in handouts mode
%   \newcommand{\Hpbreak}{\pagebreak} % pagebreaks for handouts
%   \newcommand{\Npbreak}{} % pagebreaks for printable notes
%   \newcommand{\Nonly}[1]{} % shows in printable, not in handouts
%   \newcommand{\mycleardoublepage}{\newpage} % don't need blank pages if not printing
  
%   % don't really need margins since not printing.
%   \geometry{margin=0.75in}

%   \renewcommand{\mynonumsect}[2][]{
%     \mycleardoublepage
%     \section*{#1\hskip16pt Handout: #2}
%     \addcontentsline{toc}{section}{#1 #2}
%   }

% } 



